% AY2013 制御工学 第2問（転記）
% 出典: 03_制御工学/original/03_制御工学/2013/2013se.pdf
% 要約: 伝達関数 K0, G1(s) を直列接続したシステム。ゲイン・位相、
%       ボード線図（ゲイン線図の漸近線）、ステップ応答、過渡・定常特性の考察。(50点)

\section*{第2問（50点）}

伝達関数 $K_0$, $G_1(s)$ からなる下図のシステムについて (1)〜(4) で答えよ.
なお, 入力を $U(s)$, 出力を $Y(s)$ とする.

\begin{center}
\begin{tikzpicture}[>=Latex]
  \node (in) at (0,0) {$U(s)$};
  \node (G1) [block] at (1.8,0) {$G_1(s)$};
  \node (K0) [block] at (3.8,0) {$K_0$};
  \node (out) at (5.6,0) {$Y(s)$};
  \draw[->] (in) -- (G1);
  \draw[->] (G1) -- (K0);
  \draw[->] (K0) -- (out);
\end{tikzpicture}
\end{center}

\begin{enumerate}[label=(\arabic*)]
  \item $G_1(s)=\dfrac{1}{10s^{2}+12.5s+2}$, $K_0=2$ のとき,
        システムのゲインならびに位相を求めよ.

  \item 問い (1) のボード線図のうち, ゲイン線図を漸近線によって描け.
        折れ点等の特徴点を明示し, 解答用紙の縦軸・横軸には適宜数値を記入すること.

  \item 問い (1) のパラメータのとき, システムの単位ステップ応答を求め, 概形を描け.

  \item 問い (1) と (2) の結果を用いて, このシステムの過渡特性および定常特性について
        考察し述べよ.
\end{enumerate}
